\subsubsection{The Uniform Boundedness Theorem}\label{subsec:uniform-boundedness-theorem}

It is easy to construct a sequence $f_n$ of functions on $[0,1]$ such that the sequence $f_n(x)$ is bounded for each $x$, but the $f_n$ are not uniformly bounded. For instance,
\begin{equation}\label{eq:pointwise-bounded-not-uniform-function-sequence}
f_n(x):=\begin{cases}0,&x\in[0,1/n)\\\dfrac{1}{x},&x\in[1/n,1].\end{cases}
\end{equation}
The uniform boundedness theorem says that, in contrast, for linear mappings between Banach spaces pointwise bounds imply uniform bounds.

\begin{Theorem}
\label{thm:uniform-boundedness}
Let $X,Y$ be Banach spaces. Let $T_n\in\mathcal{L}(X,Y)$, $n\in\mathbb{N}$, be such that $\|T_nx\|$ is bounded uniformly in $n$ for each fixed $x\in X$. Then $\|T_n\|$ is bounded uniformly in $n$.
\end{Theorem}

\begin{Corollary}
\label{cor:weakly-bounded-sequence-norm-bound}
Let $Z$ be a Banach space and let $x_n$ be a sequence in $Z$ such that $f(x_n)$ is bounded uniformly in $n$ for every $f\in Z^*$. Then $\|x_n\|$ is bounded uniformly in $n$.
\end{Corollary}

The proof of the corollary is obtained by setting $X=Z^*$ and $Y=\mathbb{R}$ ($\mathbb{C}$) in Theorem~\ref{thm:uniform-boundedness}.

\begin{Corollary}
\label{cor:weak-operator-bounds-imply-uniform-bound}
Let $X,Y$ be Banach spaces and let $T_n\in\mathcal{L}(X,Y)$ be such that $f(T_nx)$ is bounded uniformly in $n$ for each $f\in Y^*$ and $x\in X$. Then $\|T_n\|$ is bounded uniformly in $n$.
\end{Corollary}

For the proof, apply first Corollary~\ref{cor:weakly-bounded-sequence-norm-bound} to find that $\|T_nx\|$ is bounded uniformly in $n$ and then use Theorem~\ref{thm:uniform-boundedness}.

The proof of Theorem~\ref{thm:uniform-boundedness} is based on the Baire category theorem. We shall not carry out this proof here, but refer to the literature instead. In the following, we give a much simpler proof in a special case, namely, the case of Corollary~\ref{cor:weakly-bounded-sequence-norm-bound}, where $Z$ is a Hilbert space. Assume that the conclusion is false, i.e., $\|x_n\|$ is not uniformly bounded. Let $M=\{x_n\mid n\in\mathbb{N}\}$, and let
\begin{equation}\label{eq:hilbert-unbounded-sequence-functional-supremum}
\alpha(f)=\sup_{n\in\mathbb{N}}|(f,x_n)|
\end{equation}
for $f\in Z$. Since $M$ is not uniformly bounded, there exist $y_1\in M$ and a unit vector $e_1$ with $|(e_1,y_1)|\geq1$. Let $Q$ be the vector space spanned by $e_1$ and $y_1$. Every $x_n$ can be decomposed as $x_n=x_n^1+x_n^2$, where $x_n^1\in Q$ and $x_n^2\in Q^\perp$. Since $|(e_1,x_n)|$ and $|(y_1,x_n)|$ are uniformly bounded, the $x_n^1$ are uniformly bounded and hence the $x_n^2$ are not. We can therefore find $y_2\in M$ and a unit vector $e_2\in Q^\perp$ such that
\begin{equation}\label{eq:hilbert-unbounded-sequence-second-vector}
|(e_2,y_2)|\geq2(\alpha(e_1)+2).
\end{equation}
Inductively, we can find $y_{n+1}\in M$ and a unit vector $e_{n+1}$, orthogonal to $e_1,e_2,\ldots,e_n;y_1,y_2,\ldots,y_n$, such that
\begin{equation}\label{eq:hilbert-unbounded-sequence-inductive-vectors}
|(e_{n+1},y_{n+1})|\geq(n+1)\left(\sum_{i=1}^n\dfrac{1}{i}\alpha(e_i)+n+1\right).
\end{equation}
We now put
\begin{equation}\label{eq:hilbert-unbounded-sequence-testing-series}
f=\sum_{i=1}^{\infty}\dfrac{1}{i}e_i.
\end{equation}
This series converges because the $e_i$ are orthonormal and $\sum_{i=1}^{\infty}i^{-2}<\infty$. We obtain further
\begin{equation}\label{eq:hilbert-unbounded-sequence-contradiction}
\begin{aligned}
|(f,y_{n+1})|&=\left|\sum_{i=1}^n\dfrac{1}{i}(e_i,y_{n+1})+\dfrac{1}{n+1}(e_{n+1},y_{n+1})\right|\\
&\geq-\sum_{i=1}^n\dfrac{1}{i}\alpha(e_i)+\dfrac{1}{n+1}(n+1)\left(\sum_{i=1}^n\dfrac{1}{i}\alpha(e_i)+n+1\right)\\
&=n+1.
\end{aligned}
\end{equation}
Hence the sequence $|(f,x_n)|$ cannot be bounded, a contradiction.
